% !TeX root = ../main.tex
% ============================================
% 第七章 总结与展望
% ============================================

\chapter{总结与展望}

\section{工作总结}

本文完成的核心工作，是围绕主动健康管理场景构建一条从多模态输入组织、综合风险分析到情景推演与受控解释反馈的工程主链。针对多模态健康数据来源分散、风险分析链路不连贯以及智能问答缺乏受控约束等问题，论文完成了从理论基础、系统设计、算法建模到平台实现与实验验证的整体研究，并在统一框架下衔接了临床指标、遗传信息、文档解析结果、行为观测以及知识增强问答之间的关键关系。

从具体工作看，本文首先围绕主动健康管理平台的数据组织问题，完成了结构化临床指标、遗传信息、OCR 文档结果和行为观测信息的统一接入与管理，使后续分析能够建立在相对一致的输入基础上。在此基础上，本文设计并实现了以综合风险分析为中枢的主业务链，并在同一风险引擎基础上扩展了情景化推演能力，使平台不仅能够输出静态风险结果，还能够在当前实现边界内提供面向未来变化和干预方案的对比信息。进一步地，本文将受控 Agent 机制引入健康问答场景，通过意图分流、工具约束、证据门控、审计留痕和行为评测等机制，把问答过程限制在可追踪、可复核的运行框架内，使其承担分析结果的解释与反馈职责，而不是将其作为完全开放的生成式接口。

结合第 6 章的验证结果可以看到，当前平台已经具备较为完整的工程原型特征。训练脚本、服务逻辑和前端状态之间的主要输入边界能够保持一致，仓库级回归与浏览器场景测试也表明档案录入、OCR 状态回传、综合分析、历史会话管理和问答交互等关键链路在当前代码版本中可以稳定运行。与此同时，Agent 的受控 Harness 评测结果表明，其在固定场景下具备可回归、可审计的行为边界。总体而言，本文完成的是一套面向主动健康管理场景的工程化原型方案，其价值主要体现在把多模态风险分析、情景推演和受控问答组织成了一条可以落地验证的实现链路。

\section{不足与局限}

尽管本文已经完成多模态健康管理平台的原型构建与验证，但当前工作仍有比较明确的边界。首先，在数据与模型层面，真正高质量、长周期且完成临床表型、遗传风险和行为观测配对的真实世界数据仍然有限。当前风险建模、风险修正和情景推演更多体现为面向原型系统的工程验证，其泛化能力仍需在更大规模、更多样化的人群和场景中继续检验。尤其是第 6 章中已经说明，部分离线 benchmark 和更细粒度的消融排序并未在本轮环境下完整重跑，因此相关结论应当保持在实现基础和方法可行性的层面，而不能外推为充分的临床统计结论。

其次，在系统层面，本文实现的是一个可运行、可验证的轻量化平台原型，而不是面向大规模生产部署的终态架构。当前后端服务组织、异步任务处理、持久化方案和容器部署方式已经能够支撑论文验证与中低负载场景，但在更高并发、多用户和长时间在线运行条件下，仍可能面临资源占用、调度弹性和扩展能力方面的约束。此外，OCR 解析和问答生成均依赖外部 API 服务，其可用性、响应延迟、调用额度和隐私合规要求都会影响系统在真实环境中的稳定运行；真实运行环境下也可能先触发状态回传和待识别流程，而不是始终进入完整自动回填路径。后续若面向更严格的部署场景，还需要进一步完善本地化模型替代、请求脱敏、访问审计和服务降级策略。这些现象都说明，现有系统更适合作为平台设计思路和关键链路的验证载体，而不是直接面向复杂生产环境的最终方案。

最后，在 Agent 模块层面，本文所证明的主要是受控问答链路的工程可控性、审计可追踪性和行为一致性，而不是 Agent 在真实医疗场景中的全面临床有效性。意图分流、工具约束、证据门控和行为评测能够明显收紧系统回答边界，但这些机制并不能替代医生在复杂医学判断中的专业决策。因此，本文中的 Agent 更适合被理解为健康管理平台中的受控交互模块，而非具备独立诊疗能力的智能代理。

\section{未来展望}

后续工作首先需要继续积累更高质量的真实世界多模态健康数据，增强临床、遗传、文档与行为信息在同一主体下的长期配对能力。只有在数据覆盖度、时间跨度和标注质量进一步提升后，风险评估、风险修正和情景推演的结果才更有可能获得更稳健的统计支撑。与之对应，离线评测部分也需要进一步补足系统性的校准分析、跨人群比较和更加完整的消融实验，使模型层面的结论不再主要依赖工程侧的一致性验证。

在平台实现方面，后续工作可以在保持当前主业务链稳定的前提下，进一步增强 OCR 文档解析的鲁棒性、分析上下文的连续性以及多页面状态协同能力。对于问答模块，则更适合沿着“受控增强”而不是“能力泛化”的方向推进，在维持现有审计与约束机制的基础上，逐步提升工具利用效率、证据组织方式和人工接管协同能力。与此同时，平台若要进入更复杂的真实应用环境，还需要进一步补足隐私保护、权限治理、扩展部署和运行监测等工程能力。

总体来看，主动健康管理平台的持续发展不仅依赖模型能力增强，也依赖数据质量提升、系统工程完善和治理机制同步演进。本文完成的工作为这一方向提供了初步的工程化基础，但其更重要的意义仍然在于验证了一个可实现、可测试、可逐步扩展的技术路线。后续研究仍需要在真实应用环境中不断检验、修正并收紧其适用边界。
